\chapter{Análise de Riscos}
\label{chap:risks}

Este capítulo identifica e avalia os principais riscos
associados ao desenvolvimento e à operação da Medicano,
propondo estratégias de mitigação para cada risco
identificado.

\section{Metodologia de Análise}
\label{sec:risk-methodology}

A análise de riscos adota a abordagem qualitativa de
probabilidade~$\times$~impacto, amplamente utilizada
em gerenciamento de projetos de software
\cite{sommerville2011}. Cada risco é avaliado em duas
dimensões:

\begin{itemize}
  \item \textbf{Probabilidade}: chance de o risco se
        materializar ao longo do projeto ---
        \textit{Baixa} (B), \textit{Média} (M) ou
        \textit{Alta} (A)
  \item \textbf{Impacto}: consequência sobre o prazo,
        a qualidade ou a viabilidade do sistema caso
        o risco ocorra ---
        \textit{Baixo} (B), \textit{Médio} (M) ou
        \textit{Alto} (A)
\end{itemize}

O nível de risco resulta da combinação das duas
dimensões, conforme a Tabela~\ref{tab:risk-levels}:

\begin{table}[H]
  \centering
  \caption{Classificação do nível de risco}
  \label{tab:risk-levels}
  \begin{tabular}{lccc}
    \toprule
    \textbf{Probabilidade} $\backslash$ \textbf{Impacto}
      & \textbf{Baixo}
      & \textbf{Médio}
      & \textbf{Alto} \\
    \midrule
    Alta
      & \cellcolor{warning!30} Médio
      & \cellcolor{danger!35} Alto
      & \cellcolor{danger!35} Alto \\
    Média
      & \cellcolor{success!25} Baixo
      & \cellcolor{warning!30} Médio
      & \cellcolor{danger!35} Alto \\
    Baixa
      & \cellcolor{success!25} Baixo
      & \cellcolor{success!25} Baixo
      & \cellcolor{warning!30} Médio \\
    \bottomrule
  \end{tabular}
  \fonte{Elaborado pelos autores (2026)}
\end{table}

\section{Riscos Identificados}
\label{sec:risks-identified}

A Tabela~\ref{tab:risks} apresenta os dez riscos
identificados, organizados por categoria.

\begin{table}[H]
  \centering
  \caption{Riscos identificados e estratégias de mitigação}
  \label{tab:risks}
  \resizebox{\textwidth}{!}{%
  \begin{tabular}{p{0.6cm}p{3.2cm}p{1.8cm}p{1cm}p{1cm}p{1.2cm}p{5cm}}
    \toprule
    \textbf{ID} &
    \textbf{Descrição} &
    \textbf{Categoria} &
    \textbf{Prob.} &
    \textbf{Imp.} &
    \textbf{Nível} &
    \textbf{Estratégia de mitigação} \\
    \midrule
    R01
      & Instabilidade ou mudança de preços da API
        do modelo de linguagem
      & Técnico
      & M & A
      & \cellcolor{danger!35} Alto
      & Abstrair chamadas ao LLM em uma camada de
        serviço trocável; manter provedor alternativo
        configurado; definir orçamento mensal com
        alerta no painel do provedor \\
    \midrule
    R02
      & Alucinações do LLM gerando orientação de
        especialidade incorreta ao paciente
      & Técnico
      & M & A
      & \cellcolor{danger!35} Alto
      & Exibir aviso explícito de que o chatbot não
        substitui consulta médica; limitar respostas
        ao direcionamento de especialidade; revisar
        o \textit{system prompt} a cada sprint \\
    \midrule
    R03
      & Limitações de desempenho do cluster
        Atlas M0 em carga simultânea
      & Técnico
      & M & M
      & \cellcolor{warning!30} Médio
      & Monitorar uso via MongoDB Atlas Charts;
        implementar índices nos campos de busca;
        planejar upgrade para M2 antes do lançamento
        ao público \\
    \midrule
    R04
      & Estouro do cronograma do primeiro semestre
      & Projeto
      & M & M
      & \cellcolor{warning!30} Médio
      & Revisão semanal do burndown no Jira;
        identificar bloqueios na daily; reduzir
        escopo de funcionalidades secundárias se
        necessário \\
    \midrule
    R05
      & Expansão não controlada de escopo
        (\textit{scope creep})
      & Projeto
      & M & M
      & \cellcolor{warning!30} Médio
      & Exigir spec aprovada antes de iniciar
        qualquer sprint (SDD); registrar e priorizar
        novas demandas no backlog sem inserção
        direta no sprint ativo \\
    \midrule
    R06
      & Saída de membro da equipe durante o
        desenvolvimento
      & Projeto
      & B & A
      & \cellcolor{warning!30} Médio
      & Manter documentação atualizada no
        Confluence; adotar revisão de código por
        pares em todos os módulos críticos;
        evitar concentração de conhecimento \\
    \midrule
    R07
      & Não conformidade com a LGPD no
        tratamento de dados de saúde
      & Regulatório
      & B & A
      & \cellcolor{warning!30} Médio
      & Coletar apenas dados estritamente
        necessários; documentar o fluxo de dados
        pessoais; incluir termos de uso e política
        de privacidade antes do lançamento \\
    \midrule
    R08
      & Falha de segurança na autenticação
        ou exposição de credenciais
      & Segurança
      & B & A
      & \cellcolor{warning!30} Médio
      & Armazenar todas as credenciais no AWS
        Secrets Manager; usar Bcrypt para senhas;
        validar tokens JWT no Redis a cada
        requisição; revisar dependências com
        \texttt{npm audit} a cada sprint \\
    \midrule
    R09
      & Custos de infraestrutura AWS acima
        do orçamento previsto
      & Financeiro
      & B & M
      & \cellcolor{success!25} Baixo
      & Configurar alertas de faturamento no
        AWS Billing; utilizar instâncias do
        nível gratuito onde possível;
        revisar custos mensalmente \\
    \midrule
    R10
      & Descontinuação ou mudança de interface
        do Vercel AI SDK
      & Técnico
      & B & M
      & \cellcolor{success!25} Baixo
      & Encapsular todo o uso do SDK em um
        serviço dedicado no backend, isolando
        o restante da aplicação de mudanças
        na API do provedor \\
    \bottomrule
  \end{tabular}}
  \fonte{Elaborado pelos autores (2026)}
\end{table}

\section{Matriz de Riscos}
\label{sec:risk-matrix}

A Tabela~\ref{tab:risk-matrix} distribui os riscos
identificados na matriz de probabilidade~$\times$~impacto,
permitindo visualizar a concentração e a severidade
dos riscos do projeto.

\begin{table}[H]
  \centering
  \caption{Matriz de riscos}
  \label{tab:risk-matrix}
  \begin{tabular}{p{1.8cm}
                  >{\centering\arraybackslash}p{3.5cm}
                  >{\centering\arraybackslash}p{3.5cm}
                  >{\centering\arraybackslash}p{3.5cm}}
    \toprule
    &
    \textbf{Impacto Baixo} &
    \textbf{Impacto Médio} &
    \textbf{Impacto Alto} \\
    \midrule
    \textbf{Prob. Alta}
      & \cellcolor{warning!30} ---
      & \cellcolor{danger!35} ---
      & \cellcolor{danger!35} --- \\[18pt]
    \textbf{Prob. Média}
      & \cellcolor{success!25} ---
      & \cellcolor{warning!30} R03 \quad R04 \quad R05
      & \cellcolor{danger!35} R01 \quad R02 \\[18pt]
    \textbf{Prob. Baixa}
      & \cellcolor{success!25} ---
      & \cellcolor{success!25} R09 \quad R10
      & \cellcolor{warning!30} R06 \quad R07 \quad R08 \\[18pt]
    \bottomrule
  \end{tabular}
  \fonte{Elaborado pelos autores (2026)}
\end{table}

Os dois riscos de nível alto (R01 e R02) estão ambos
relacionados ao componente de triagem por LLM, que
representa o principal diferencial técnico da
plataforma e, ao mesmo tempo, sua maior dependência
externa. As mitigações adotadas — abstração do
provedor e avisos ao usuário — reduzem o impacto
desses riscos sem comprometer a funcionalidade
central do sistema.
